\documentclass[10pt,twocolumn]{article}
\usepackage[margin=0.75in]{geometry}
\usepackage{mathptmx,microtype,amsmath,booktabs}
\usepackage[numbers,sort&compress]{natbib}
\usepackage{url,hyperref}
\hypersetup{colorlinks=true,urlcolor=blue,citecolor=blue,linkcolor=blue,
  pdftitle={Adaptive Image Reads for Agent Training},
  pdfauthor={Majd Zayyad}}
\setlength{\emergencystretch}{1.5em}
\Urlmuskip=0mu plus 1mu
\title{\bfseries Adaptive Image Reads for Agent Training}
\author{Majd Zayyad}
\date{November 2026}
\begin{document}
\maketitle

\begin{abstract}
An agent-training job can leave its learner waiting for tools that are
still fetching files from a remote image. Reading larger chunks reduces
requests but can waste bandwidth; smaller chunks have the opposite cost.
We study a reader that chooses how much to fetch from recent access
patterns and observed request cost. The question is whether this local
decision shortens training execution outside model inference.
We combine measured image-read jobs with a scheduler that counts setup,
tool computation, queues, verification, teardown, and learner updates,
while excluding inference service and its queue.
\input{generated/abstract-result}
The experiments explain when image-read savings reach the learner, and
when other work absorbs them.
\end{abstract}

\section{Why image reads matter}
Agentic training repeatedly turns model output into actions: inspect a
repository, import a library, compile a program, or run tests. Those tools
need files inside their sandbox. With remote-backed images, some of the
files arrive only when read. If tools finish late, a training update can
wait even when the learner itself is ready.

DSec makes the workload mismatch concrete. It reuses 3FS, its
training-storage backend, for on-demand sandbox images. Section~5.3
reports strong throughput for large sequential I/O but poor performance
for small random requests~\cite{dsec2026}. DSec bridges that gap with
local writes, local metadata where possible, and bulk reads into caches.
Its OverlayBD/ublk path fetches 256\,KiB chunks. Agent tools, however,
often want a few files promptly, then move to another short burst of work.

This follows the familiar \emph{I/O-granularity trade-off}.
Large fetches amortize request overhead; small ones limit read
amplification---bytes transferred but not used. Lambda's image-loading
study describes this same random-access versus sequential-throughput
trade-off~\cite{lambda2023}. The opportunity is to adapt the client-side
fetch size to the current access pattern, rather than require a new
backend or a larger cache.

Read-ahead already adapts to sequential access, as in AMP~\cite{amp2007}.
TIP and Khameleon connect prefetching decisions to cost or expected
usefulness~\cite{tip1995,khameleon2020}. AdaCache selects aligned blocks
for known missing intervals~\cite{adacache2023}; WaRP plans object-store
requests from intervals already known to be required~\cite{warp2026}.
Our reader instead chooses a blocking refill before the next tool read
is known. It combines likely continuation with the cost of the aligned
windows that the current read would cross.

We therefore ask: \emph{does adaptive fetching reduce the time to finish
the same training-execution work, excluding model inference?}
This isolates the execution stack from model-serving latency without
reducing the amount of tool or learner work.

\input{generated/results-table}

\section{Reader and training experiment}
\paragraph{Choosing a refill.}
The reader considers aligned windows from 16\,KiB to 1\,MiB. Recent
adjacent reads suggest that a stream will continue; jumps reduce that
expectation. Completed requests update its estimate of fixed request
cost and transfer rate. It selects the window with the lowest
\[
  \frac{\text{estimated fetch time}}{\text{bytes likely to be used}}.
\]
The cost includes crossing multiple windows, not just the requested
byte count. Cache capacity, replacement, workers, and in-flight limits
are unchanged. Extra bytes arrive inside the blocking refill; there is
no asynchronous prefetcher. The implementation and parameter definitions
are in the accompanying code.

\paragraph{Measured tool-stage input.}
An independent C++ reader accesses a 512\,MiB image over loopback HTTP
with a 64\,MiB cache. Random, sequential, and mixed traces are tested at three
concurrency levels and under three controlled latency/bandwidth settings, giving
27 conditions. Each has three measured repetitions. Timing covers the
entire replay process, including setup and drain. Only client-observable
timings reach the controller.

All controls use the same executable. Alongside fixed256, we test other
fixed sizes, choose one static setting on separate development traces,
and include a simple grow/reset rule. That rule increases its window on
successive missed demands and resets after a jump. Every measurement
checks the returned data and common resource budgets.

\paragraph{Execution time excluding inference.}
We use each measured replay duration as the cost of one synthetic
image-read job, not as a pure storage latency. A rollout consists of
setup, the image-read job, additional tool computation, verification,
and teardown. Model-call placeholders take zero time and bypass the
model-worker queue. We resimulate the schedule; subtracting summed
inference durations from overlapping work would give the wrong clock.
Actor, CPU, and tool-worker limits remain unchanged.
A learner update starts only after its assigned rollouts finish; the
next update's rollouts wait for that update.

Every simulated run completes the same 256 rollouts and eight learner
updates. Policies share job order, service-time variation, non-inference
work, and resources. We vary tool computation and learner-update cost
to expose different reasons execution might wait.
The clock stops after the final update, not after the last image read.

Each tool worker has an independent reader/backend budget, and each
image-read job starts with an empty reader cache. This is a controlled
execution-time model: it does not simulate learning quality, real
repository tasks, warm sandbox reuse, or contention inside a shared
3FS service. Its clock is inference-excluded execution, not actual
end-to-end training wall time. Complete parameters and individual
outcomes are provided with the data.

\section{What changes execution time?}
\input{generated/results-text}

Image fetching deserves attention when it delays usable rollouts.
Even without inference, CPU work and learner updates can limit the gain.
Correct stream identity remains essential. Captured coding-tool traces
and matched training runs are needed to assess actual time-to-quality.

\paragraph{Code and reproduction.}
Source, calibration data, the full simulation matrix, and commands for
recollecting the measurements and rebuilding this paper are public at
\url{https://github.com/majdzayyad/agent-training-image-io}.
Release \texttt{v1.1.0} fixes the version used here.

\bibliographystyle{plainnat}
\bibliography{background,refill}
\end{document}
